\section*{Hoofdstuk \ref{ch:b3:structural-analysis} --- Structurele analyse}

\begin{solution}{exo:b3:structural-analysis:1}
De \omterm{def:b3:structural-analysis:stereo}{pool} heeft richting $40 + 180 = \ang{220}$ en een duik van $90 - 30 = \ang{60}$. Het
vlak helt met $90 - 20 = \ang{70}$ naar $315 - 180 = \ang{135}$: $135/70$.
\end{solution}

\begin{solution}{exo:b3:structural-analysis:2}
$\theta = \ang{60}$. \omterm{def:b3:structural-analysis:stereo}{Hoekgetrouw net}: $10\tan\ang{30} = \qty{5.77}{cm}$;
oppervlaktegetrouw net: $10\sqrt2\sin\ang{30} = \qty{7.07}{cm}$.
\end{solution}

\begin{solution}{exo:b3:structural-analysis:3}
$\vec n_1 = (\sin\ang{30}, 0, \cos\ang{30})$, $\vec n_2 = (-\sin\ang{30}, 0,
\cos\ang{30})$, $\vec n_1\cdot\vec n_2 = -0.25 + 0.75 = 0.5$: de polen liggen
\ang{60} uit elkaar, de interflankhoek is \ang{120}, een open plooi.\par
\end{solution}

\begin{solution}{exo:b3:structural-analysis:4}
$\vec d = (\cos\ang{60}, 0, -\sin\ang{60}) = (0.5, 0, -0.866)$ en $\vec u =
0.707\,\vec s + 0.707\,\vec d = (0.354, 0.707, -0.612)$: richting
$\operatorname{atan2}(0.354, 0.707) = \ang{26.6}$, duik $\arcsin0.612 =
\ang{37.8}$. De rake ligt tussen \ang{-90} (normaal) en 0 (sinistraal): een
schuine afschuiving met een sinistrale component.
\end{solution}

\begin{solution}{exo:b3:structural-analysis:5}
\begin{gather*}
  \vec n_1 = (0.5567, 0.3214, 0.7660),\qquad \vec n_2 = (-0.4698, -0.1710, 0.8660),\\
  \vec n_1\times\vec n_2 = (0.4093, -0.8420, 0.0558),
\end{gather*}
\sloppy met lengte $0.9379$. Het tegengestelde, $-\vec n_1\times\vec n_2$, wijst omlaag:
richting \ang{334.1}, duik $\arcsin(0.0558/0.9379) = \ang{3.4}$. $\vec n_1\cdot\vec
n_2 = 0.3469$, \ang{69.7} tussen de normalen, de interflankhoek \ang{110.3}:
een open plooi.\par
\end{solution}

\begin{solution}{exo:b3:structural-analysis:6}
Een duik boven \ang{60} betekent $\theta < \ang{30}$. Oppervlaktegetrouw net: $r^2 =
2\sin^2\ang{15} = 0.134$, \qty{13.4}{\percent} van het net; \omterm{def:b3:structural-analysis:stereo}{hoekgetrouw net}:
$\tan^2\ang{15} = 0.072$, \qty{7.2}{\percent}. Op de bol bevat de kap
$1 - \cos\ang{30} = 0.134$ van het halfrond: alleen het oppervlaktegetrouwe net toont
het met zijn werkelijke aandeel, daarom worden polen erop geteld.
\end{solution}

\begin{solution}{exo:b3:structural-analysis:7}
\begin{gather*}
  T_{xx} = \frac{2\sin^2\ang{20}}3 = 0.078,\qquad
  T_{zz} = \frac{2\cos^2\ang{20} + 1}3 = 0.922,\\
  T_{xz} = -\frac{\sin\ang{20}\cos\ang{20} - \sin\ang{20}\cos\ang{20}}3 = 0,
\end{gather*}
en elke term met $y$ is 0. De eigenwaarden zijn $0.922$ (verticaal),
$0.078$ (oost) en $0$ (noord): de \omterm{def:b3:structural-analysis:pi}{bèta-as} is horizontaal, noord--zuid.
De drie polen liggen in één vlak, dus geen ervan heeft een component langs zijn normaal,
en $\vec u^{\mathsf T}T\vec u = 0$ voor $\vec u$ langs die normaal.
\end{solution}

\begin{solution}{exo:b3:structural-analysis:8}
$a = 2$, $b = 1.5$, $k = 1/0.5 = 2$: insnoering. $Z = 1/(1.8\times1.1) =
0.505$, $a = 1.636$, $b = 2.178$, $k = 0.636/1.178 = 0.54$: afplatting.
\end{solution}

\begin{solution}{exo:b3:structural-analysis:9}
$S_1/S_2 = 2.25$ en $S_1S_2 = 1$: $S_1 = 1.5$, $S_2 = 0.667$; de lange as
werd met \qty{50}{\percent} verlengd. De holte heeft de halve assen $1.5\times0.4 =
\qty{0.60}{mm}$ en $0.667\times0.4 = \qty{0.27}{mm}$.
\end{solution}

\begin{solution}{exo:b3:structural-analysis:10}
\omterm{def:b3:structural-analysis:stereo}{Hoekgetrouw net}: $1/(2\cos^2\ang{30}) = 0.667$ langs de straal en
$\tan\ang{30}/\sin\ang{60} = 0.577/0.866 = 0.667$ in de omtrekrichting: gelijk, conform.
Oppervlaktegetrouw net: $\mathrm dr/(R\,\mathrm d\theta) = (\sqrt2/2)\cos\ang{30} = 0.612$
en $\sqrt2\sin\ang{30}/\sin\ang{60} = 0.707/0.866 = 0.816$; het product is
$0.500$, de oppervlaktefactor van de stelling. Een kleine cirkel wordt een ovaal met assen
in de verhouding $0.816/0.612 = 1.33$, met de lange as rond het middelpunt.
\end{solution}

\begin{solution}{exo:b3:structural-analysis:11}
De breuk is verticaal, dus de normaal ervan is horizontaal en loodrecht op de
verticale $\sigma_2$; volgens deel (b) van de stelling staat de gliding loodrecht op
$\sigma_2$, horizontaal, voor elke $\Phi$. Met $\sigma_1 = 1$ in noordelijke richting,
$\sigma_3 = 0$ in oostelijke richting: $S\vec n = (0, -0.342, 0)$, $\vec n\cdot S\vec n =
0.117$, $\vec\tau = (0, -0.342, 0) - 0.117\,(0.940, -0.342, 0) = (-0.110,
-0.302, 0)$ en $-\vec\tau = (0.110, 0.302, 0)$, naar \ang{020}. Het blok
aan de kant van $\vec n$, de oostkant, beweegt noord-noordoost langs de
strekking: gezien vanaf de westkant beweegt het naar links, een sinistrale breuk (rake
0). $\sigma_2$ speelde in de berekening geen rol.
\end{solution}

\begin{solution}{exo:b3:structural-analysis:12}
$S = (7.5 - 6.0)/1.0 = \qty{1.5}{km}$; $h = 6.0/1.5 = \qty{4.0}{km}$ onder de
onderste horizon. De doorsnede was $20 + 1.5 = \qty{21.5}{km}$ lang en veranderde
met $-1.5/21.5 = \qty{-7.0}{\percent}$.
\end{solution}

\begin{solution}{pb:b3:structural-analysis:1}
\textbf{1.} De \omterm{def:b3:structural-analysis:stereo}{pool} heeft richting \ang{255} en een duik van \ang{77}; $\theta =
\ang{13}$ en $r = 10\sqrt2\sin\ang{6.5} = \qty{1.60}{cm}$ vanaf het middelpunt,
naar \ang{255}.

\textbf{2.} $\vec n = (\sin\ang{13}\sin\ang{75}, \sin\ang{13}\cos\ang{75},
\cos\ang{13}) = (0.2173, 0.0582, 0.9744)$. Zes bijna evenwijdige eenheidsvectoren
tellen op tot bijna 6: $|\vec S_E| = 5.9974$.

\textbf{3.} Oostflank: $\bar\alpha = \operatorname{atan2}(1.2583, 0.2017) =
\ang{80.9}$ en $\bar\delta = \arccos(5.8604/5.9974) = \ang{12.3}$; $\bar R
= 5.9974/6 = 0.99956$ en de spreiding is $\sqrt{2\times0.00044} =
\qty{0.0297}{rad}$, of \ang{1.7}. Westflank: $\bar\alpha =
\operatorname{atan2}(-1.2445, 0.1886) = \ang{278.6}$ en $\bar\delta =
\arccos(5.8636/5.9972) = \ang{12.1}$; $\bar R = 0.99953$, spreiding \ang{1.8}.

\textbf{4.}
\[
  \vec S_E\times\vec S_W = (0.0774, -14.6714, 0.4883),
\]
met lengte $14.680$. Het tegengestelde wijst naar beneden: de richting is $\operatorname{atan2}(-0.0774, 14.6714) = \ang{359.7}$ en de
duik $\arcsin(0.4883/14.680) = \ang{1.9}$.

\textbf{5.} De cosinus van de hoek tussen de gemiddelde normalen is
$32.835/35.968 = 0.9129$, het inproduct van de sommen gedeeld door hun
lengtes: \ang{24.1} tussen de normalen, een interflankhoek van \ang{155.9}:
een open plooi, dicht bij zacht.

\textbf{6.}
\begin{gather*}
  \vec m_E = (0.2098, 0.0336, 0.9772),\qquad \vec m_W = (-0.2075, 0.0314, 0.9777),\\
  \vec m_E - \vec m_W = (0.4173, 0.0022, -0.0005):
\end{gather*}
een horizontale normaal naar \ang{089.7}. Het assenvlak is verticaal, met strekking
\ang{359.7}, praktisch noord--zuid: de plooi is rechtopstaand.

\textbf{7.} De gordel is het vlak waarvan de \omterm{def:b3:structural-analysis:stereo}{pool} de \omterm{def:b3:structural-analysis:pi}{bèta-as} is: hellingrichting
\ang{179.7}, helling $90 - 1.9 = \ang{88.1}$, bijna het oost--west verticale
vlak. Elke flank bevat de plooias, dus elke normaal staat loodrecht
op die as en ligt in het vlak dat daar loodrecht op staat.

\textbf{8.} $0.04400 + 0.00144 + 0.95456 = 1.00000$. $\vec b^{\mathsf T}T\vec
b = 0.00038$, de kleinste eigenwaarde: $\vec b$ is, binnen de afronding, de
eigenvector $\vec e_3$, en de gordel door de twee gemiddelde polen is de
best passende gordel van alle twaalf.

\textbf{9.} $\arcsin\sqrt{0.00038} = \arcsin0.0195 = \ang{1.1}$.

\textbf{10.} Langs de gordel liggen de polen op ongeveer \ang{12.2} van de verticaal,
aan weerszijden; $\sin^2\ang{12.2} = 0.045$, het gemiddelde kwadraat van hun
oost--west-componenten, wat $\lambda_2$ is.

\textbf{11.} Ja: de polen wijken gemiddeld \ang{1.1} van de gordel af, niet
meer dan de spreiding van \ang{1.7} van de metingen van één flank rond hun gemiddelde.

\textbf{12.} $\ang{1.1}/\sqrt{12} = \ang{0.3}$: de duik is \ang{1.9} $\pm$
\ang{0.3}.

\textbf{13.} Een rake van \ang{-90}: een zuivere afschuiving, dip-slip. Volgens Anderson:
$\sigma_1$ verticaal, $\sigma_3$ horizontaal langs de hellingrichting,
\ang{120}--\ang{300}, $\sigma_2$ horizontaal langs de strekking,
\ang{030}--\ang{210}.

\textbf{14.} De normaal van $120/60$ ligt in het verticale vlak met richting
\ang{120}, dat $\vec e_1$ en $\vec e_3$ bevat: $n_2 = 0$. Volgens onderdeel (b) is de
gliding loodrecht op $\sigma_2$, de strekking, voor elke $\Phi$: het is
dip-slip. De breuk zegt niets over $\Phi$; in de formule verdwijnen zowel $n_2$ als
$\vec e_2\cdot\vec w$, en daarmee zowel teller als noemer.

\textbf{15.} $\vec n = (\sin\ang{55}\sin\ang{150}, \sin\ang{55}\cos\ang{150},
\cos\ang{55}) = (0.4096, -0.7094, 0.5736)$; $n_1 = 0.574$; $n_2 = 0.5\times
0.4096 + 0.866\times(-0.7094) = -0.410$.

\textbf{16.} $\vec e_1\cdot\vec w = 0.198$ en $\vec e_2\cdot\vec w = 0.5\times
0.771 + 0.866\times0.605 = 0.909$, dus $\Phi = -(0.574\times0.198)/(-0.410\times
0.909) = 0.1137/0.3727 = 0.305$.

\textbf{17.} $(0.305 + 0.294 + 0.301 + 0.306)/4 = 0.30$. De rakes zijn
afgerond op de graad; één graad rake verschuift $\Phi$ op deze breuken met ongeveer 0.02,
dus de afrondingsfout van een halve graad verklaart de spreiding.

\textbf{18.} $\sigma_2 - \sigma_3 = 0.3(\sigma_1 - \sigma_3)$: $\sigma_2$ ligt
dichter bij $\sigma_3$ dan bij $\sigma_1$. Beide horizontale spanningen liggen ruim
onder de verticale; het gebied werd vooral naar \ang{120} uitgerekt, een
beetje langs \ang{030}.

\textbf{19.} Met $\vec n = (0.750, -0.433, 0.500)$ en de gliding $\vec u =
(0.433, -0.250, -0.866)$, $(\vec n - \vec u)/\sqrt2$ duikt onder \ang{75} richting
\ang{300}: P; $(\vec n + \vec u)/\sqrt2$ duikt \ang{15} richting \ang{120}:
T. Ze liggen in hetzelfde verticale vlak als $\sigma_1$ en $\sigma_3$, maar
\ang{15} daar vandaan, omdat de breuk \ang{30}, niet \ang{45}, maakt met
$\sigma_1$.

\textbf{20.} $\tan a = \tan\ang{12.3}\times\cos(\ang{90} - \ang{80.9}) =
0.2180\times0.9874 = 0.2153$: $a = \ang{12.1}$.

\textbf{21.} $L_0 = 8000/\cos\ang{12} = \qty{8178.7}{m}$, $S =
\qty{178.7}{m}$, een relatieve verandering van $\cos\ang{12} - 1 =
\qty{-2.2}{\percent}$.

\textbf{22.} $4000\tan\ang{12} = \qty{850.2}{m}$; $A = \tfrac12\times8000\times
850.2 = \qty{3.40e6}{m^2}$.

\textbf{23.} $h = \num{3.40e6}/178.7 = \qty{19.0}{km}$. Een décollement aan de
basis van de bedekking zou $h$ van enkele honderden meters geven en een oppervlakteoverschot
ongeveer vijftig keer kleiner: de plooien zijn daar niet losgekoppeld. Ofwel plooide het
fundament samen met de bedekking, of de synclinalen zakten onder het oorspronkelijke
niveau, wat $A$ kleiner maakt; alleen de doorsnede kan dat niet uitmaken.

\textbf{24.} Horizontale verplaatsing $200/\tan\ang{60} = \qty{115.5}{m}$, gliding $200/\sin\ang{60}
= \qty{230.9}{m}$.

\textbf{25.} $\tan a = \tan\ang{60}\cos\ang{30} = 1.500$: $a = \ang{56.3}$, en
de schijnbare horizontale verplaatsing is $200/1.5 = \qty{133.3}{m}$. De gliding van de breuk is
omlaag langs de helling, richting \ang{120}: de doorsnede langs \ang{120} toont ware
lengten en is dus de doorsnede die hersteld moet worden.

\textbf{26.} De breuk snijdt de plooien door, dus is zij jonger en wordt als eerste hersteld,
waarbij de hangende wand over \qty{230.9}{m} terug omhoog langs de breuk schuift; daarna worden de plooien
ontplooid. De \omterm{def:b3:structural-analysis:restoration}{pinlijnen} liggen op de assenvlakken van de synclinalen, in de
liggende wand, waar de symmetrische plooien geen gliding tussen de lagen hadden.

\textbf{27.} De plooias van Varden duikt
\textbf{\ang[text-series-to-math=true]{1.9}} naar het noorden (richting
\ang{359.7}): de \ang{2} die in het gebied is ingebouwd, teruggevonden tot op de
\ang{0.3} die twaalf metingen toelaten.
\end{solution}
